% First complete-results writing pass, for independent scientific review.
% Source: paper/array_revision_fair/analysis/full; no new inference or fits.
% Replace the Results body, not the Methods. Table fragments are complete tables.
% Place the existing validated Figure 2 at the insertion marker below.
% Figure 3 below references the existing full-analysis PDF; do not regenerate it.

\subsection{Complete evaluation}
Both stages completed all 57 datasets and five paired outer partitions,
with no missing or skipped task. The fixed stage contains 1,710
depth--feature blocks, 5,130 single-tree scores, and 136,800 forest
scores. The selected stage adds 855 separately refitted forests, three
families for each of the 285 outer tasks. Table~\ref{tab:performance}
reports representative fixed architectures and all selected families.
Differences below are percentage points of accuracy. Intervals are
pointwise, unadjusted $95\%$ paired dataset-bootstrap intervals for
the mean effect. Table~\ref{tab:comparisons} gives the eleven primary
contrasts with one joint Holm adjustment; other comparisons are descriptive.
Scores and costs first average repetitions within datasets, then weight
datasets equally.

\subsection{All-feature single trees and pure forests}
At $D=3$, mean accuracies for single trees with $k=1,2,3$ are
$67.82\%,70.59\%,71.36\%$. The paired horizon-two/three gains over
CART are $2.77$ points [$0.68$, $5.22$] and $3.54$ points [$1.45$,
$6.01$], but medians are only $0.13$ and $0.40$ points, with
$29/3/25$ and $31/3/23$ wins/ties/losses. Mean fit-only times are
$0.00103$, $0.0282$, and $1.684$ s, respectively. At $D=6$, mean
differences shrink to $0.39/0.59$ points and medians become
$-0.48/-1.27$ points. Without a depth limit, mean differences are
$-5.13$ [$-8.69$, $-1.84$] and $-8.29$ [$-12.45$, $-4.39$]
points. Mean realized depths then increase from $10.1$ for CART to
$15.1/17.0$ for horizons two/three. Larger realized trees and poorer
test scores coincide here, without identifying their cause.

For pure $T=100$ forests with all features, horizon-two/three mean
differences from the matched CART forest are $1.90/2.34$ points at
$D=3$, $-0.58/-1.09$ at $D=6$, and $-6.84/-9.68$ without a depth
limit. Shallow mean fitting work is $0.0677$, $1.6807$, and $80.2307$
s for the three pure horizons. With square-root feature subsampling,
the corresponding mean differences are $2.38/2.85$, $1.27/1.96$,
and $1.08/0.78$ points. The last two intervals are [$0.35$, $1.88$]
and [$-0.75$, $2.44$], with medians $0.29$ and zero points.
These are contrasts with each feature policy's own CART control,
not a formal interaction test or evidence that subsampling universally
improves absolute accuracy. The unrestricted all-feature losses must
not be extrapolated to the square-root-feature setting.

\subsection{Sparse horizon replacements}
At $T=100$, $D=3$, with all features, replacing ten CART members with
horizon-two or horizon-three members gives mean gains of $1.06$ points
[$0.45$, $1.77$] and $1.27$ points [$0.61$, $2.02$]. Both medians
are zero; wins/ties/losses are $27/21/9$ and $28/14/15$.
Their joint-Holm signed-rank values are $p_H=0.00996$ and $0.02066$.
Balanced-accuracy gains are $1.15/1.33$ points. Mean selected-slot
fitting work rises from $0.0677$ s for CART to $0.2293/8.0933$ s,
or $3.39/119.52$ times the control. These are architecture-matched,
not equal-computation gains. The descriptive $85/10/5\%$ composition
has a mean gain of $1.36$ points [$0.52$, $2.30$] at $4.2260$ s.

With square-root features, the two shallow replacement gains are
$0.42$ [$0.07$, $0.78$] and $1.04$ [$0.46$, $1.74$] points,
with $p_H=0.09539$ and $0.00150$. Their medians are zero and
$0.13$ points. Without a depth limit, all four primary replacement
means are small: $0.05/-0.02$ points with all features and
$0.22/0.28$ points with square-root features. All four medians are
zero; adjusted values range from $0.34732$ to $1$. This is not
equivalence evidence. At the descriptive depth-six all-feature
setting, the replacement mean differences are $0.27/0.58$ points,
again with zero medians.

For the two shallow all-feature anchors, equal weighting of the 48
known dataset families reduces the means to $0.77/0.91$ points,
with intervals [$0.25$, $1.45$] and [$0.32$, $1.63$]. On the 38
tasks outside the named synthetic/constructed set, the means are
$0.53/0.63$ points, with intervals [$0.06$, $1.22$] and [$0.07$,
$1.35$]. These sensitivities concern these anchors and do not make
the previously explored convenience sample independent or representative.

\subsection{Descriptive mean-cost curves}
% PARENT FIGURE2 INSERTION: put the validated Figure 2 environment here.
% BEGIN FIGURE2_POSTPLOT_TEXT: retain directly after that environment.
Figure~\ref{fig:tradeoff} shows all sixteen depth-three, all-feature
compositions at $T=20,40,60,100,200$. Neither tree count nor the
farther-horizon fraction gives uniformly monotone observed mean
accuracy. Among plotted points with mean fitting work at most $0.7$
s, the pure horizon-two $T=20$ forest has the highest observed mean:
$74.79\%$ at $0.3381$ s. A $T=100$ forest with $25\%$ horizon-two
members has $74.57\%$ accuracy at $0.4709$ s. Sparse mixtures
therefore do not necessarily improve on smaller pure farther-sighted
forests. They offer other tradeoffs: the $T=20$, $25\%$ horizon-two
forest has $74.06\%$ accuracy at $0.0955$ s. Every plotted composition
containing horizon-three members has mean fitting work above $0.7$ s.
This reference is a retrospective comparison of benchmark means, not
a selected budget policy or a per-dataset feasibility guarantee.
Connecting lines do not represent evaluated intermediate models or
certified equal-cost comparisons, and the figure contains no paired
uncertainty. Its fitting-work axis excludes preparation, prediction,
and selection cost.
% END FIGURE2_POSTPLOT_TEXT

\subsection{Training-selected families}
% New figure environment: paths are relative to paper/array_revision/main.tex.
\begin{figure*}[t]
\centering
\includegraphics[width=\textwidth]{../array_revision_fair/analysis/full/Fig3_tuned_families.pdf}
\caption{Training-only selected families on the complete 57-dataset sample and five paired repetitions. Family 1/2 allows horizons one and two; Family 1/2/3 allows all three. Both can select pure forests. Left: mean paired accuracy differences with pointwise, unadjusted $95\%$ dataset-bootstrap intervals, not intervals for the Holm-adjusted signed-rank tests. Right: standalone cross-validation constituent fitting work, direct selected-model refit wall time, and the sum of measured workflow stage wall times. The last measure excludes checkpoint I/O, startup, downloads, and idle gaps; it is not an uninterrupted end-to-end stopwatch. Common validation-bank work is fully charged to every family requiring it.}
\label{fig:tuned}
\end{figure*}

Mean accuracy is $78.99\%$ for the RF family, $79.96\%$ for the
one-/two-sighted family, and $80.56\%$ for the inclusive family
(Figure~\ref{fig:tuned}). Relative to selected RF, the one-/two-sighted
family has a mean gain of $0.97$ points [$-0.07$, $2.30$], median
$0.00617$ points, $29/8/20$ wins/ties/losses, and $p_H=1$.
The inclusive family gains $1.57$ points [$0.30$, $3.28$], with
median zero, $28/11/18$ wins/ties/losses, and $p_H=0.48838$.
Its balanced-accuracy gain is $1.60$ points. Inclusive versus
one-/two-sighted gives $0.60$ points [$-0.11$, $1.63$], median
zero, and $p_H=1$. The positive inclusive mean interval and the
non-rejected signed-rank location null concern different estimands,
not contradictory decisions about one hypothesis. Neither supports
calling the methods equivalent or saying that no mean improvement
was observed.

Selection does not always choose a mixed model. Among the 285
inclusive choices, 133 are pure forests ($61/28/44$ at horizons
one/two/three) and 152 are mixed (146 use two horizons and six use all three).
The one-/two-sighted family selects 168 pure forests and 117
mixtures. Thus the aggregate selected-family gain cannot be
attributed solely to sparse replacement or heterogeneity.
Depth-three/six/unrestricted selection counts are $66/87/132$ for RF,
$90/106/89$ for the one-/two-sighted family, and $104/103/78$ for
the inclusive family. Both all-feature and square-root-feature
settings are selected in every family. The RF label therefore
denotes a library-estimator candidate space including all-feature
bagged CART, not exclusively square-root-feature ensembles.

\subsection{Selection costs and sensitivity of the selected-family effects}
Standalone cross-validation fitting work averages $2.271$, $112.079$,
and $2,671.049$ s for RF, one-/two-sighted, and inclusive families.
These charges include all banks required by each search space,
even when a pure CART forest is ultimately selected. Direct refit
wall times, including construction/bootstrap preparation but not
test prediction, are $0.0626$, $0.9518$, and $12.9374$ s. Their
ratios to RF are $15.20$ and $206.61$ for the two expanded families;
these are refit ratios, not selection-cost ratios. Mean sums of
measured workflow stage wall times are $3.405$, $118.035$, and
$2,693.803$ s. They are not uninterrupted end-to-end times.
Shared banks are fully attributed to each standalone family but
counted only once in the separate physical accounting. Fixed-forest
constituent work and selected-model direct refit wall time have
different boundaries and are not interchangeable.

The inclusive-versus-RF mean effect becomes $1.03$ points
[$-0.27$, $2.90$] with equal weighting of the 48 known families
and $0.53$ points [$-0.91$, $2.77$] on the 38 tasks outside the named
synthetic/constructed set. Both intervals cross zero. The corresponding one-/two-sighted
effects are $0.61$ [$-0.32$, $2.02$] and $0.38$ [$-0.66$, $2.09$]
points. These prespecified analyses show that the estimated tuned
mean benefits are sensitive to benchmark weighting and membership;
they do not establish an absence of benefit or a broadly transferable
accuracy--cost advantage.
